\chapter{Εισαγωγή}

Ο τουρισμός ήταν για δεκαετίες αμιγώς θετικό φαινόμενο. Αποτελούσε
καθαρή πηγή εσόδων, ανάπτυξης και πολιτισμικής ανταλλαγής. Μέχρι
και τα τέλη του 20ου αιώνα, το πρόβλημα ήταν πάντα πώς να
προσελκύσεις περισσότερους τουρίστες. Η στροφή ήρθε ξαφνικά, χάρη
στη ραγδαία ανάπτυξη της τεχνολογίας, της ψηφιοποίησης και
γενικότερα την αύξηση του βιοτικού επιπέδου παγκοσμίως. Οι
εταιρίες ενοικίασης καταλυμάτων, τα φθηνά αεροπορικά εισιτήρια,
τα μέσα κοινωνικής δικτύωσης και η παγκοσμιοποίηση συμβάλλουν στη
δημιουργία μαζικών τουριστικών ροών σε πόλεις που δεν είναι
σχεδιασμένες να τις απορροφήσουν. Ξαφνικά πόλεις όπως η Βενετία,
νησιά όπως η Σαντορίνη, αλλά ακόμα και μικρά χωριά βρέθηκαν να
$``$πνίγονται$"$ από επισκέπτες. Ο όρος υπερτουρισμός
\en{(overtourism)} εμφανίστηκε στη βιβλιογραφία γύρω στο 2017
και είναι το φαινόμενο που καλείται να αναλύσει αυτή η εργασία.

\section{Σημασία του προβλήματος}

Ο τουρισμός από τη φύση του έχει πολλαπλά θετικά οφέλη. Πρώτο και
κύριο όφελος αποτελεί τον οικονομικό αντίκτυπο που αφήνει στην
τοπική κοινωνία. Δημιουργούνται νέες θέσεις εργασίας, αυξάνεται το
ΑΕΠ, ενισχύονται οι τοπικές επιχειρήσεις και προσελκύονται νέες
επενδύσεις. Επιπλέον, ενισχύονται οι υποδομές, όπως το οδικό δίκτυο,
οι μεταφορές και οι δημόσιοι χώροι, από τους οποίους επωφελούνται
και οι μόνιμοι κάτοικοι. Τέλος, με τον τουρισμό προωθείται η
πολιτιστική ανταλλαγή, η αλληλοκατανόηση και η συνύπαρξη
διαφορετικών κουλτούρων.

Ωστόσο, το βασικό ζήτημα προκύπτει όταν ο τουρισμός δεν είναι πλέον
διαχειρίσιμος, καθώς οι αρνητικές επιπτώσεις του υπερτουρισμού
ποικίλλουν και επηρεάζουν διαφορετικές ομάδες με διαφορετικό τρόπο.
Για τους κατοίκους, οι επιπτώσεις εκδηλώνονται κυρίως ως υποβάθμιση
της ποιότητας ζωής, αυξημένος θόρυβος, κυκλοφοριακή συμφόρηση, άνοδος
των τιμών ενοικίων και σταδιακός εκτοπισμός από τις κεντρικές γειτονιές
της πόλης. Το περιβάλλον, από την άλλη, δέχεται έντονη ρύπανση, φθορά
μνημείων και υπερφόρτωση των υποδομών ύδρευσης, αποκομιδής απορριμμάτων
και ενέργειας. Τέλος, ακόμα και οι ίδιοι οι τουρίστες δεν παραμένουν
ανεπηρέαστοι, αφού η υπερπλήρωση των αξιοθέατων αυξάνει τους χρόνους
αναμονής, διογκώνει το κόστος του ταξιδιού και υποβαθμίζει συνολικά
την τουριστική εμπειρία.
Επομένως, η μελέτη και η διαχείριση τουριστικών ροών είναι ζωτικής
σημασίας για τις ανάγκες ευημερίας των σύγχρονων κοινωνιών.

Η διαχείριση του υπερτουρισμού απαιτεί κατανόηση του \emph{πώς} και
\emph{πού} κινούνται οι τουρίστες, ποιες περιοχές υπερφορτώνονται
και πώς επηρεάζονται οι μόνιμοι κάτοικοι. Αυτό δεν είναι εύκολο
να μελετηθεί αποκλειστικά με παραδοσιακές στατιστικές μεθόδους.
Χρειαζόμαστε προσομοίωση που να αναπαράγει τη δυναμική του
φαινομένου σε πραγματικό αστικό περιβάλλον.

\section{Στόχοι της Διπλωματικής Εργασίας}

Ο κύριος στόχος είναι η ανάπτυξη ενός \en{agent-based model (ABM)}
που να προσομοιώνει τον υπερτουρισμό σε πραγματικό αστικό χώρο.
Πιο συγκεκριμένα, η διπλωματική εργασία ξεκινά με την κατασκευή ενός
ρεαλιστικού οδικού δικτύου βασισμένου σε ελεύθερα δεδομένα από το
\en{OpenStreetMap}, πάνω στο οποίο μοντελοποιούνται δύο τύποι
πρακτόρων: οι μόνιμοι κάτοικοι και οι τουρίστες. Στη συνέχεια,
εξομοιώνεται η έξυπνη κίνηση των τουριστών εντός της πόλης,
καθώς αυτοί επισκέπτονται αξιοθέατα, εστιατόρια και καταστήματα,
και μετράται η επίδραση αυτής της τουριστικής δραστηριότητας τόσο
στην ευτυχία των κατοίκων όσο και στην ικανοποίηση των ίδιων των
τουριστών. Τέλος, τα αποτελέσματα οπτικοποιούνται σε πραγματικό χάρτη,
με τη χρήση \en{heatmaps} θορύβου και ρύπανσης.

\section{Συνεισφορά της Διπλωματικής Εργασίας}

Η συγκεκριμένη δουλειά προσφέρει πέντε βασικά καινοτόμα στοιχεία σε
σχέση με υπάρχουσες προσεγγίσεις. Πρώτον, η κατασκευή του χώρου
βασίζεται σε πραγματικά δεδομένα \en{GIS} και όχι σε τυχαίο πλέγμα
\en{(grid)}, εξασφαλίζοντας ρεαλιστική αναπαράσταση του αστικού
περιβάλλοντος. Δεύτερον, εφαρμόζεται υβριδική εύρεση διαδρομών που
συνδυάζει την κατεύθυνση προς τον στόχο με την ελκυστικότητα του δρόμου,
προσομοιώνοντας πιο πιστά τη συμπεριφορά πεζών τουριστών. Τρίτον, τα
σημεία ενδιαφέροντος ιεραρχούνται σε τρία επίπεδα με διαφορετική
βαρύτητα στη συμπεριφορά των πρακτόρων, κάτι που δεν συναντάται συχνά
σε αντίστοιχα μοντέλα. Τέταρτον, παρέχεται πλήρης \en{web-based}
οπτικοποίηση σε πραγματικό χάρτη, διευκολύνοντας την παρακολούθηση της
προσομοίωσης σε πραγματικό χρόνο. Τέλος, το μοντέλο διαθέτει μεταφερσιμότητα,
καθώς μπορεί να εφαρμοστεί σε οποιαδήποτε πόλη ή χωριό της Ελλάδας.

\section{Διάρθρωση της Διπλωματικής Εργασίας}

Η εργασία αυτή είναι οργανωμένη σε επτά κεφάλαια: Στο Κεφάλαιο 2
δίνεται το θεωρητικό υπόβαθρο των βασικών τεχνολογιών που
σχετίζονται με τη διπλωματική αυτή. Αρχικά περιγράφονται τα
\en{Agent Based Models (ABMs)}, στη συνέχεια αναλύονται τα
χαρακτηριστικά των γεωχωρικών δεδομένων και οι βασικές έννοιες
της θεωρίας γράφων και τέλος παραθέτονται οι τεχνολογίες που
χρησιμοποιήθηκαν για την υλοποίηση. Στο
Κεφάλαιο 3 αρχικά περιγράφονται οι σχετικές με το θέμα εργασίες
και στη συνέχεια δίνεται ο στόχος της συγκεκριμένης εργασίας. Στο
Κεφάλαιο 4 παρουσιάζεται η ανάλυση και η σχεδίαση του συστήματος,
δηλαδή η περιγραφή των υποσυστημάτων και των εφαρμογών του. Η
περιγραφή της υλοποίησης του συστήματος, με ανάλυση των βασικών
αλγορίθμων καθώς και λεπτομέρειες σχετικά με τις πλατφόρμες και τα
προγραμματιστικά εργαλεία που χρησιμοποιήθηκαν δίνεται στο
Κεφάλαιο 5. Στο Κεφάλαιο 6 παρουσιάζονται τα αποτελέσματα και
η αξιολόγηση πειραμάτων προσομοίωσης για ποικίλα σενάρια, καθώς
και τα συμπεράσματα που προκύπτουν από αυτή τη μελέτη.
Τέλος στο Κεφάλαιο 7 δίνεται η συνεισφορά αυτής της
διπλωματικής εργασίας, καθώς και μελλοντικές επεκτάσεις.